% ch07 译文：原文 §5a（main.tex 第 1805–2169 行）
\section{$G$-等变层\label{sec:G-equiv-sheaves}}
本节回顾 $G$-等变层的定义（可参看 \cite[\href{https://stacks.math.columbia.edu/tag/03LE}{Tag 03LE}]{stacks-project}），并建立概形 $X$ 上 $G$-等变凝聚层与向量丛同 $X$ 上 $H$-等变凝聚层与向量丛之间的关系，其中 $H\subseteq G$ 为子群。

\subsection{$G$-等变 $\mathcal{O}_X$-模的范畴}

设 $G=(G, m, e)$ 为 $\field$ 上的群概形，乘法为 $m : G \times G \rightarrow G$，单位截面为 $e : \Spec(\field) \rightarrow G$；设 $X$ 为 $\field$ 上的概形，带有 $G$-作用 $\sigma : G \times X \rightarrow X$。若干基本定义见 \cite[\href{https://stacks.math.columbia.edu/tag/022R}{Section 022R}]{stacks-project} 与 \cite[\href{https://stacks.math.columbia.edu/tag/01II}{Section 01II}]{stacks-project}。

\begin{defn}\label{defn:EquivariantSheaves}
一个\emph{$G$-等变 $\mathcal{O}_X$-模}
$(\mathcal{F}, \theta)$
指一个 $\mathcal{O}_X$-模 $\mathcal{F}$，连同 $\mathcal{O}_{G\times X}$-模的同构
$$
\theta : \sigma^* \mathcal{F} \xrightarrow{\cong} \left(\pr_X^{G \times X}\right)^* \mathcal{F}
$$
（称为 $\mathcal{F}$ 的\emph{$G$-线性化}），其中 $\pr_X^{G \times X} : G \times X \rightarrow X$ 为自然投影，并且下列条件成立：
% [原文 p.31]
\begin{itemize}
\item \emph{（余循环条件）}
$
(m \times \id_X)^* \theta
=
(\id_G \times \sigma)^* \theta \comp \left(\pr_{G \times X}^{G \times G \times X}\right)^* \theta.
$

\item \emph{（单位条件）}
$(e \times \id_X)^*\theta = \id_{\mathcal{F}}$.
\end{itemize}

从\emph{$G$-等变 $\mathcal{O}_X$-模}
$\mathcal{F}=(\mathcal{F}, \theta)$
到
$\mathcal{F}'=(\mathcal{F}', \theta')$
的\emph{态射}，指一个 $\mathcal{O}_X$-模态射
$
f : \mathcal{F} \rightarrow \mathcal{F}'
$
，它与 $G$-线性化交换，即
$$
\sigma^* f \comp \theta'
=
\theta \comp \left(\pr_X^{G \times X}\right)^* f
$$
作为 $\mathcal{O}_{G \times X}$-模的态射成立。

我们把\emph{$G$-等变 $\mathcal{O}_X$-模的范畴}记作
$\Mod_G(X)$。
若 $(\mathcal{F},\theta)$ 底下的 $\mathcal{F}$ 作为 $\mathcal{O}_X$-模具有相应性质，就说 $G$-等变 $\mathcal{O}_X$-模 $(\mathcal{F},\theta)$ 是
\emph{拟凝聚}、\emph{凝聚}
或\emph{有限秩局部自由}的。以 $G$-等变拟凝聚、凝聚、有限秩局部自由的 $\mathcal{O}_X$-模为对象，$\Mod_G(X)$ 有三个满子范畴，分别记作
$\QCoh_G(X)$、
$\Coh_G(X)$、
$\VB_G(X)$，
它们满足
$$
\VB_G(X)\subset\Coh_G(X)\subset\QCoh_G(X)\subset\Mod_G(X).
$$
\end{defn}

\begin{conceptbox}{$G$-等变层}
$G$-等变层（\emph{$G$-equivariant sheaf}）就是叠加了「沿群作用的函子性」的层：$G$-线性化 $\theta:\sigma^*\mathcal{F}\xrightarrow{\cong}\left(\pr_X^{G\times X}\right)^*\mathcal{F}$ 把「先作用再取层」与「先取层再作用」两种拉回等同起来，余循环条件保证该等同在 $G\times G\times X$ 上相容，单位条件则让单位截面处的等同退化为恒等。它等价于商叠 $[X/G]$ 上的拟凝聚层，因而是「带 $G$-作用的空间上的层」的正确范畴化：普通层只记得 $X$ 的局部结构，等变层额外记住每点沿 $G$-轨道的平移。本文需要它，是因为 B 侧的 $\mathbb{G}_m$-主丛 $L(Y,\xi)^{\times}\to Y$ 上的层正是 $Y$ 上的 $\mathbb{G}_m$-等变层，而等变层按权分解（见 §5c）恰好与 A 侧绕 $\mathbb{P}^1$ 纤维方向的单值性、停止位选择相对应。
\par\smallskip\noindent{\footnotesize 参考：nLab，\emph{equivariant sheaf}；Stacks Project，Tag 03LE。}
\end{conceptbox}

\begin{yizhu}[$\Theta$ 的共轭方向]
内部 Hom 的线性化取共轭 $\Theta:\alpha\mapsto\theta^{-1}\comp\alpha\comp\theta'$，两头方向相反：源侧用 $\theta^{-1}$ 把 $\sigma^*\mathcal{F}$ 拉回 $\left(\pr_X^{G\times X}\right)^*\mathcal{F}$，靶侧用 $\theta'$ 把 $\sigma^*\mathcal{F}'$ 送到 $\left(\pr_X^{G\times X}\right)^*\mathcal{F}'$。若误写成 $\theta\comp\alpha\comp\theta'^{-1}$，得到的共轭与原意相反，§5c 中权 $\cha$ 会整体变号。还要注意本节证明里出现的 $\cdot$ 有时表示复合而非数乘，如下面命题证明中的 $\theta_U^{-1}\cdot\sigma_U^*f\cdot\theta'_U$。
\end{yizhu}

\begin{defn}[等变层的内部张量积]
\label{defn:TensorProductOfEquivariantModules}
设 $(\mathcal{F}, \theta)$ 与 $(\mathcal{F}', \theta')$ 为 $\Mod_G(X)$ 中的两个 $G$-等变 $\mathcal{O}_X$-模。
$(\mathcal{F}, \theta)$ 与 $(\mathcal{F}', \theta')$ 的\emph{（内部）张量积}
$$
(\mathcal{F}, \theta)
\underset{\mathcal{O}_X}{\otimes}
(\mathcal{F}', \theta')
:=
\left(\mathcal{F}\underset{\mathcal{O}_X}{\otimes}\mathcal{F}', \, \theta\otimes\theta'\right)
\in
\Mod_G(X),
$$
定义为 $\mathcal{O}_X$-模 $\mathcal{F}\underset{\mathcal{O}_X}{\otimes}\mathcal{F}'$
连同 $G$-线性化
$$
\theta\otimes\theta'
:
\sigma^*\left(
\mathcal{F}
\underset{\mathcal{O}_X}{\otimes}
\mathcal{F}'
\right)
\xrightarrow{\cong}
\left(\pr_X^{G \times X}\right)^*
\left(\mathcal{F}
\underset{\mathcal{O}_X}{\otimes}
\mathcal{F}'\right)
$$
给出，后者是张量积 $\theta \otimes \theta'$ 在典范同构
$$
\begin{gathered}
\sigma^*\left(\mathcal{F}\underset{\mathcal{O}_X}{\otimes}\mathcal{F}'\right)
\cong
\sigma^* \mathcal{F} \underset{\mathcal{O}_{G \times X}}{\otimes} \sigma^* \mathcal{F}' \\
\text{与}\quad
\left(\pr_X^{G \times X}\right)^* \left(\mathcal{F}\underset{\mathcal{O}_X}{\otimes}\mathcal{F}'\right)
\cong
\left(\pr_X^{G \times X}\right)^*\mathcal{F} \underset{\mathcal{O}_{G \times X}}{\otimes} \left(\pr_X^{G \times X}\right)^*\mathcal{F}'
\end{gathered}
$$
之下的像。

张量积的\emph{单位对象}定义为
$(\mathcal{O}_X, \id_{\mathcal{O}_{G\times X}}) \in \Mod_G(X)$，
其中
$
\id_{\mathcal{O}_{G\times X}}
:
\sigma^* \mathcal{O}_X
\xrightarrow{\cong}
\left(\pr_X^{G \times X}\right)^* \mathcal{O}_X
$
由层的典范同构
$
\sigma^* \mathcal{O}_X
\cong
\mathcal{O}_{G\times X}
$
与
$
\left(\pr_X^{G \times X}\right)^* \mathcal{O}_X
\cong
\mathcal{O}_{G\times X}
$
诱导得到。

配备该张量积，
$\Mod_G(X)$
构成一个
\emph{对称幺半范畴}
$
\left(
\Mod_G(X),
\underset{\mathcal{O}_X}{\otimes},
\left(\mathcal{O}_X,\theta(0)\right)
\right).
$
此外，$\underset{\mathcal{O}_X}{\otimes}$ 限制下去，也赋予
$\QCoh_G(X)$、$\Coh_G(X)$、$\VB_G(X)$
对称幺半范畴结构。

\end{defn}

\begin{yizhu}[单位对象为何写成 $\theta(0)$]
本定义的开头把张量积的单位对象写成 $(\mathcal{O}_X,\id_{\mathcal{O}_{G\times X}})$，末尾的对称幺半范畴却写成 $\left(\mathcal{O}_X,\theta(0)\right)$：两处指的是同一个线性化，$\theta(0)$ 应理解为「未用任何特征扭过的恒等线性化」，即上面由 $\sigma^*\mathcal{O}_X\cong\mathcal{O}_{G\times X}\cong\left(\pr_X^{G\times X}\right)^*\mathcal{O}_X$ 诱导的那一个。原文记号不统一，读时按平凡线性化处理即可；§5b 的扭构造正是在这个 $\theta(0)$ 上乘以特征。
\end{yizhu}

\begin{defn}[等变层的内部 Hom]
\label{defn:InternalHomOfEquivariantModules}
设 $(\mathcal{F}, \theta)$ 与 $(\mathcal{F}', \theta')$ 为 $\Mod_G(X)$ 中的两个 $G$-等变 $\mathcal{O}_X$-模，
且 $\mathcal{F}$ 有限表现\footnote{
$\mathcal{O}_X$-模 $\mathcal{F}\in\Mod(X)$
称为有限表现，若
(1)
$\mathcal{F}\in\VB(X)$ 为有限秩局部自由，
或
(2)
$X$ 局部 Noether 且 $\mathcal{F}\in\Coh(X)$ 为凝聚层。
}。
$(\mathcal{F}, \theta)$ 与 $(\mathcal{F}', \theta')$ 的\emph{内部 Hom}
$$
\mathcal{H}om_{\mathcal{O}_X}\left( (\mathcal{F}, \theta), (\mathcal{F}', \theta') \right)
:=
\left( \mathcal{H}om_{\mathcal{O}_X}(\mathcal{F}, \mathcal{F}'), \, \Theta \right)
\in
\Mod_G(X),
$$
定义为 $\mathcal{O}_X$-模 $\mathcal{H}om_{\mathcal{O}_X}(\mathcal{F}, \mathcal{F}')$
连同 $G$-线性化
$$
\Theta
:
\sigma^*\mathcal{H}om_{\mathcal{O}_X}(\mathcal{F}, \mathcal{F}')
\xrightarrow{\cong}
\left(\pr_X^{G \times X}\right)^* \mathcal{H}om_{\mathcal{O}_X}(\mathcal{F}, \mathcal{F}')
$$
给出，后者由共轭
$
\alpha
\mapsto
\theta^{-1}
\comp
\alpha
\comp
\theta'
$
在典范同构\footnote{
这些典范同态是同构，因为态射
$\sigma$、
$\pr_X^{G\times X}
:
G \times X \rightarrow X$
都平坦（任何群概形 $G$ 对基平坦），
且 $\mathcal{F}$ 有限表现。
}
之下给出
% [原文 p.32]
\begin{align*}
\sigma^* \mathcal{H}om_{\mathcal{O}_X}(\mathcal{F}, \mathcal{F}')
&\xrightarrow{\cong}
\mathcal{H}om_{\mathcal{O}_{G \times X}}\left(\sigma^* \mathcal{F}, \sigma^* \mathcal{F}'\right)
\quad\text{与}\quad
\\[2mm]
\left(\pr_X^{G \times X}\right)^* \mathcal{H}om_{\mathcal{O}_X}(\mathcal{F}, \mathcal{F}')
&\xrightarrow{\cong}
\mathcal{H}om_{\mathcal{O}_{G \times X}}\left(\left(\pr_X^{G \times X}\right)^*\mathcal{F}, \left(\pr_X^{G \times X}\right)^*\mathcal{F}'\right).
\end{align*}

更确切地说，对每个开集 $U \subseteq G \times X$ 与局部截面
$
\alpha \in \left(\sigma^*\mathcal{H}om_{\mathcal{O}_X}(\mathcal{F}, \mathcal{F}')\right)(U)
$，
把它看作 $\mathcal{O}_{G\times X}|_U$-模同态
$
\alpha : \sigma^*\mathcal{F}|_U \rightarrow \sigma^*\mathcal{F}'|_U
$，
截面 $\Theta(U)(\alpha)$ 由复合
$$
\left(\theta|_U\right)^{-1}
\comp
\alpha
\comp
\theta'|_U
:
\left(\pr_X^{G \times X}\right)^*\mathcal{F}|_U
\rightarrow
\left(\pr_X^{G \times X}\right)^*\mathcal{F}'|_U
$$
定义。

配备张量积与内部 Hom 后，
$\VB_G(X)$
（以及当 $X$ 局部 Noether 时的 $\Coh_G(X)$）
构成一个\emph{闭对称幺半范畴}。

\end{defn}

\begin{yizhu}[内部 Hom 为何要假设有限表现]
把 $\mathcal{H}om$ 搬过拉回需要两个条件：$\sigma$ 与 $\pr_X^{G\times X}$ 平坦（群概形对基平坦，这一条自动成立），以及 $\mathcal{F}$ 有限表现。缺了后者，$\sigma^*\mathcal{H}om(\mathcal{F},\mathcal{F}')$ 与 $\mathcal{H}om(\sigma^*\mathcal{F},\sigma^*\mathcal{F}')$ 只是同态而非同构，$\Theta$ 便无法定义成线性化。原文用两条脚注把这两处补上；读时注意「有限表现」在局部 Noether 情形下就是凝聚，而在一般概形上则要单独假设。
\end{yizhu}

本段最后，我们计算 $G$-等变层之间的截面。
\begin{prop}\label{prop:InvariantSectionsOfEquivariantInternalHom}
设
$(\mathcal H,\Theta)
:=
\mathcal Hom_{\mathcal O_X}
\bigl((\mathcal F,\theta),(\mathcal F',\theta')\bigr)$。
则对每个 $G$-不变开子概形 $U\subseteq X$，有
$$
\mathcal H(U)^G
=
\Hom_{\Mod_G(U)}
\left(
(\mathcal F,\theta)|_U,\,
(\mathcal F',\theta')|_U
\right).
$$
\end{prop}

\begin{proof}
截面 $f\in\mathcal H(U)$ 是 $\mathcal O_U$-模态射 $f:\mathcal F|_U\rightarrow\mathcal F'|_U$。
由 $G$-不变截面的定义与 $\Theta$ 的定义，下列等价成立：
\begin{align*}
f\in\mathcal H(U)^G
&\quad\Leftrightarrow\quad
\Theta|_{G\times U}(\sigma_U^*f)
=
\left(\pr_U^{G\times U}\right)^*f
\\
&\quad\Leftrightarrow\quad
\theta_U^{-1}\cdot\sigma_U^*f\cdot\theta'_U
=
\left(\pr_U^{G\times U}\right)^*f
\\
&\quad\Leftrightarrow\quad
\sigma_U^*f\cdot\theta'_U
=
\theta_U\cdot
\left(\pr_U^{G\times U}\right)^*f
\end{align*}
最后一个等式恰是 $f$ 的 $G$-等变性条件。
\end{proof}
